\setcounter{section}{1}
\section{CÔNG THỨC LƯỢNG GIÁC}
\subsection{LÝ THUYẾT CẦN NHỚ}
\subsubsection{Công thức cộng}
\begin{khung4}{Công thức cộng}
	\begin{multicols}{2}
		\begin{itemize}
			\item $\cos (a-b) = \cos a \cos b + \sin a\sin b$.
			\item $\cos (a+b) = \cos a \cos b - \sin a\sin b$.
			\item $\sin (a-b) = \sin a \cos b - \sin b \cos a$.
			\item $\sin (a+b) = \sin a \cos b + \sin b \cos a$.
			\item $\tan (a-b) = \dfrac{\tan a - \tan b}{1+\tan a \tan b}$.
			\item $\tan (a+b) = \dfrac{\tan a + \tan b}{1-\tan a \tan b}$.
		\end{itemize}
	\end{multicols}
\end{khung4}
\begin{khung4}{Trường hợp đặc biệt}
	\begin{itemize}
		\item $\sin x + \cos x = \sqrt{2}\sin\left(x+\dfrac{\pi}{4}\right) = \sqrt{2}\cos \left(x-\dfrac{\pi}{4}\right)$.
		\item $\sqrt{3}\sin x + \cos x = 2\sin \left(x+\dfrac{\pi}{6}\right) = 2\cos \left(x-\dfrac{\pi}{3}\right)$.
		\item $\sin x + \sqrt{3}\cos x = 2\sin \left(x+ \dfrac{\pi}{3}\right) = 2\cos \left(x-\dfrac{\pi}{6}\right)$.
	\end{itemize}
\end{khung4}
\subsubsection{Công thức nhân đôi}
\begin{multicols}{2}
	\begin{khung4}{Công thức nhân đôi}
		\begin{itemize}
			\item $\sin 2a = 2\sin a \cos a$.
			\item $\cos 2a = \cos^2a-\sin^2a \\= 2\cos^2a-1 = 1-2\sin^2a$.
			\item $\tan 2a = \dfrac{2\tan a}{1-\tan^2a}$.
		\end{itemize}
	\end{khung4}
	\begin{khung4}{Công thức hạ bậc}
		\begin{itemize}
			\item $\sin^2a= \dfrac{1-\cos 2a}{2}$.
			\item $\cos^2a = \dfrac{1+\cos 2a}{2}$.
			\item $\tan^2a=\dfrac{1-\cos2a}{1+\cos 2a}$.
		\end{itemize}
	\end{khung4}
\end{multicols}

% \begin{note}
% 	Áp dụng công thức cộng cho $3a = a +2a$, ta có công thức nhân ba:
% \end{note}
\begin{khung4}{Công thức nhân ba}
	\begin{multicols}{2}
		\begin{itemize}
			\item $\sin3a= 3\sin a -4\sin^3a$.
			\item $\cos3a= 4\cos^3a-3\cos a$.
			\item $\tan3a = \dfrac{3\tan a - \tan^3 a}{1-3\tan^2a}$.
		\end{itemize}
	\end{multicols}
\end{khung4}

\subsubsection{Công thức biến đổi tích thành tổng}
\begin{khung4}{Công thức tích thành tổng}
	\begin{multicols}{2}
		\begin{itemize}
			\item $\cos a \cos b = \dfrac{1}{2}\left[\cos (a-b) + \cos (a+b)\right]$.
			\item $\sin a \sin b = \dfrac{1}{2}\left[\cos (a-b)-\cos(a+b)\right]$.
			\item $\sin a \cos b = \dfrac{1}{2}\left[\sin (a-b)+\sin (a+b)\right]$.
		\end{itemize}
	\end{multicols}
\end{khung4}
\subsubsection{Công thức biến đổi tổng thành tích}
% Công thức biến đổi tổng thành tích được xây dựng bằng cách $a=\dfrac{a+b}{2}, b = \dfrac{a-b}{2}$ trong công thức biến đổi tích thành tổng.
\begin{khung4}{Công thức tổng thành tích}
	\begin{multicols}{2}
		\begin{itemize}
			\item $\cos a+ \cos b = 2\cos\dfrac{a+b}{2}\cos \dfrac{a-b}{2}$.
			\item $\cos a- \cos b = -2\sin\dfrac{a+b}{2}\sin \dfrac{a-b}{2}$.
			\item $\sin a+ \sin b = 2\sin\dfrac{a+b}{2}\cos \dfrac{a-b}{2}$.
			\item $\sin a -\sin b = 2\cos\dfrac{a+b}{2}\sin \dfrac{a-b}{2}$.
		\end{itemize}
	\end{multicols}
\end{khung4}
\subsection{PHÂN LOẠI VÀ PHƯƠNG PHÁP GIẢI TOÁN}
\begin{dang}{Sử dụng công thức cộng, công thức nhân đôi}
\end{dang}

\begin{vd}
	Khai triển:
	\begin{tasks}(1)
		\task $\cos \left(x+\dfrac{\pi}{3}\right)$\dotfill;
		\task $\cos \left(x-\dfrac{\pi}{6}\right)$\dotfill;
		\task $\sin \left(x+60^{\circ}\right)$\dotfill;
		\task $\sin \left(\dfrac{\pi}{6}-x\right)$\dotfill;
		\task $\tan \left(a+\dfrac{\pi}{3}\right)$\dotfill;
		\task $\tan \left(\dfrac{\pi}{6}-2 a\right)$\dotfill;
	\end{tasks}
	\loigiai{
		\begin{enumerate}[a)]
			\item $\cos\left(x + \dfrac{\pi}{3}\right) = \cos x \cos\dfrac{\pi}{3} - \sin x \sin\dfrac{\pi}{3} = \dfrac{1}{2}\cos x - \dfrac{\sqrt{3}}{2}\sin x$.
			\item  $\cos\left(x - \dfrac{\pi}{6}\right) = \cos x \cos\dfrac{\pi}{6} + \sin x \sin\dfrac{\pi}{6} = \dfrac{\sqrt{3}}{2}\cos x + \dfrac{1}{2}\sin x$.
			\item $\sin\left(x + 60^\circ\right) = \sin x \cos 60^\circ + \cos x \sin 60^\circ = \dfrac{1}{2}\sin x + \dfrac{\sqrt{3}}{2}\cos x$.
			\item $\sin\left(\dfrac{\pi}{6} - x\right) = \sin\dfrac{\pi}{6} \cos x - \cos\dfrac{\pi}{6} \sin x = \dfrac{1}{2}\cos x - \dfrac{\sqrt{3}}{2}\sin x$.
			\item $\tan\left(a + \dfrac{\pi}{3}\right) = \dfrac{\tan a + \tan\dfrac{\pi}{3}}{1 - \tan a \tan\dfrac{\pi}{3}} = \dfrac{\tan a + \sqrt{3}}{1 - \sqrt{3}\tan a}$.
			\item $\tan\left(\dfrac{\pi}{6} - 2a\right) = \dfrac{\tan\dfrac{\pi}{6} - \tan 2a}{1 + \tan\dfrac{\pi}{6} \tan 2a} = \dfrac{\dfrac{\sqrt{3}}{3} - \tan 2a}{1 + \dfrac{\sqrt{3}}{3} \tan 2a} = \dfrac{ - \sqrt{3}\tan 2a}{\sqrt{3} + \tan 2a}$.
		\end{enumerate}
	}
\end{vd}
% \dongcham{22}
\begin{vd}
	Chứng minh các đẳng thức sau:
	\begin{tasks}(1)
		\task $\cos a+\sin a=\sqrt{2}\cos \left(\dfrac{\pi}{4}-a\right)=\sqrt{2}\sin \left(\dfrac{\pi}{4}+a\right)$.
		\task $\cos a-\sin a=\sqrt{2}\cos \left(\dfrac{\pi}{4}+a\right)=\sqrt{2}\sin \left(\dfrac{\pi}{4}-a\right)$.
		\task $\sqrt{3}\sin x + \cos x = 2\sin \left(x+\dfrac{\pi}{6}\right) = 2\cos \left(x-\dfrac{\pi}{3}\right)$.
		\task $\sin x + \sqrt{3}\cos x = 2\sin \left(x+ \dfrac{\pi}{3}\right) = 2\cos \left(x-\dfrac{\pi}{6}\right)$.
	\end{tasks}
	\loigiai{
		\begin{enumerate}[a)]
			\item Ta có $\cos a+\sin a = \sqrt{2}\left(\dfrac{1}{\sqrt{2}}\cos a+\dfrac{1}{\sqrt{2}}\sin a\right)=\sqrt{2}\left(\cos a\cos\dfrac{\pi}{4}+\sin a\sin\dfrac{\pi}{4}\right)=\sqrt{2}\cos\left(\dfrac{\pi}{4}-a\right)$.\\
			      Mặt khác ta cũng có $\sqrt{2}\left(\dfrac{1}{\sqrt{2}}\cos a+\dfrac{1}{\sqrt{2}}\sin a\right)=\sqrt{2}\left(\sin\dfrac{\pi}{4}\cos a+\cos\dfrac{\pi}{4}\sin 2\right)=\sqrt{2}\sin\left(\dfrac{\pi}{4}+a\right)$.
			\item Ta có $\cos a-\sin a = \sqrt{2}\left(\dfrac{1}{\sqrt{2}}\cos a-\dfrac{1}{\sqrt{2}}\sin a\right)=\sqrt{2}\left(\cos a\cos\dfrac{\pi}{4}-\sin a\sin\dfrac{\pi}{4}\right)=\sqrt{2}\cos\left(\dfrac{\pi}{4}+a\right)$.\\
			      Mặt khác ta cũng có $\sqrt{2}\left(\dfrac{1}{\sqrt{2}}\cos a-\dfrac{1}{\sqrt{2}}\sin a\right)=\sqrt{2}\left(\sin\dfrac{\pi}{4}\cos a-\cos\dfrac{\pi}{4}\sin 2\right)=\sqrt{2}\sin\left(\dfrac{\pi}{4}-a\right)$.
			\item Ta có $\sqrt{3}\sin x + \cos x = 2 \left( \dfrac{\sqrt{3}}{2} \sin x + \dfrac{1}{2} \cos x \right) = 2 \left( \sin x \cos \dfrac{\pi}{6} + \cos x \sin \dfrac{\pi}{6}\right) = 2\sin \left(x+\dfrac{\pi}{6}\right) $.\\
			Mặt khác ta cũng có $2 \left( \dfrac{\sqrt{3}}{2} \sin x + \dfrac{1}{2} \cos x \right) = 2 \left( \sin x \sin \dfrac{\pi}{3} + \cos x \cos \dfrac{\pi}{3}\right) = 2\cos \left(x-\dfrac{\pi}{3}\right) $.
			\item Ta có $\sin x + \sqrt{3}\cos x = 2 \left( \dfrac{1}{2} \sin x + \dfrac{\sqrt{3}}{2} \cos x \right) = 2 \left( \sin x \cos \dfrac{\pi}{3} + \cos x \sin \dfrac{\pi}{3}\right) = 2\sin \left(x+\dfrac{\pi}{3}\right) $.\\
			Mặt khác ta cũng có $2 \left( \dfrac{1}{2} \sin x + \dfrac{\sqrt{3}}{2} \cos x \right) = 2 \left( \sin x \sin \dfrac{\pi}{6} + \cos x \cos \dfrac{\pi}{6}\right) = 2\cos \left(x-\dfrac{\pi}{6}\right) $.
		\end{enumerate}
	}
\end{vd} \cham{13}
\begin{vd}
	Cho $\sin\alpha=-\dfrac{3}{5}$ và $\dfrac{3\pi}{2}<\alpha<2\pi.$ Tính $\cos\alpha$,  $\tan\alpha$; $\cos2\alpha$ và $\sin\left( \alpha+\dfrac{19\pi}{4}\right)$.
	\loigiai{
		Ta có $\cos^2\alpha = 1-\sin^2 \alpha = 1-\dfrac{9}{25}=\dfrac{16}{25} \Rightarrow \cos \alpha = \pm \dfrac{4}{5}$. \\
		Do $\dfrac{3\pi}{2}<\alpha<2\pi$ nên $\cos \alpha >0$, suy ra $\cos \alpha = \dfrac{4}{5}$.
		\begin{itemize}
			\item [$\bullet$] $\tan\alpha =\dfrac{\sin\alpha}{\cos \alpha}=-\dfrac{3}{4}$.
			\item [$\bullet$] $\cos 2\alpha = 1-2\sin^2\alpha = \dfrac{7}{25}$.
			\item [$\bullet$] $\sin\left(\alpha +\dfrac{19\pi}{4}\right)=\sin\alpha \cdot\cos \left(\dfrac{19\pi}{4}\right)+\cos\alpha \cdot\sin \left(\dfrac{19\pi}{4}\right)=\dfrac{7\sqrt{2}}{10}$.
		\end{itemize}
	}
\end{vd} \dongcham{4}

\begin{vd}
	Cho $\tan\alpha=-2$ và $\dfrac{\pi}{2}<\alpha<\pi.$ Tính $\cos\alpha$, $\cos\left( \alpha-\dfrac{3\pi}{4}\right)$ và $\tan \left( \alpha+\dfrac{\pi}{4}\right) $.
	\loigiai{
		Ta có $\cos^2 \alpha =\dfrac{1}{1+\tan^2 \alpha}=\dfrac{1}{5} \Rightarrow \cos \alpha = \pm \dfrac{\sqrt{5}}{5}$.\\
		Do $\dfrac{\pi}{2}<\alpha<\pi$ nên $\cos \alpha <0$, suy ra $\cos \alpha = - \dfrac{\sqrt{5}}{5}$.
		\begin{itemize}
			\item [$\bullet$] $\sin \alpha = \tan \alpha \cdot \cos \alpha =\dfrac{2\sqrt{5}}{5}$.
			\item [$\bullet$] $\cos\left(\alpha -\dfrac{3\pi}{4}\right)=\cos\alpha \cdot\cos \left(\dfrac{3\pi}{4}\right)+\sin\alpha \cdot\sin \left(\dfrac{3\pi}{4}\right)=\dfrac{3\sqrt{10}}{10}$.
			\item [$\bullet$] $\tan \left( \alpha+\dfrac{\pi}{4}\right)=\dfrac{\tan \alpha + \tan\dfrac{\pi}{4}}{1-\tan \alpha \cdot \tan \dfrac{\pi}{4}}=-\dfrac{1}{3}$.
		\end{itemize}}
\end{vd} \dongcham{4}
\begin{vd}%[TH]%[DCHT Toán 11 - KNTT -Phong Lê] %[1K1B2-1] 
	Tính giá trị của biểu thức $P = \dfrac{\cos\dfrac{5\pi}{18}\cos\dfrac{\pi}{9}+\sin \dfrac{5\pi}{18}\sin \dfrac{\pi}{9} }{\sin \dfrac{\pi}{5} \cos \dfrac{3\pi}{10} + \cos \dfrac{\pi}{5}\sin\dfrac{3\pi}{10}}$.
	\dapso{$P = \dfrac{\sqrt{3}}{2}$.}
	\loigiai{Ta có $P = \dfrac{\cos \left(\dfrac{5\pi}{18}-\dfrac{\pi}{9}\right)}{\sin\left(\dfrac{\pi}{5}+\dfrac{3\pi}{10}\right)} = \dfrac{\cos \dfrac{\pi}{6}}{\sin \dfrac{\pi}{2}} = \dfrac{\sqrt{3}}{2}$.\\
		Vậy $P = \dfrac{\sqrt{3}}{2}$.}
\end{vd}
\begin{vd}%[VDT]%[DCHT Toán 11 - KNTT -Phong Lê] %[1K1K2-1] 
	Chứng minh giá trị của biểu thức
	$$P = \sin \left(\dfrac{\pi}{6}-\alpha\right) + \sin \left(\dfrac{\pi}{6} + \alpha\right) -\cos \alpha$$
	không phụ thuộc vào $\alpha$.
	% \dapso{$P=0$.}
	\loigiai{
		Ta có\\
		$\begin{array}{ll}
				P & = \sin \left(\dfrac{\pi}{6}-\alpha\right) + \sin \left(\dfrac{\pi}{6} + \alpha\right) -\cos \alpha                                              \\
				  & =\sin \dfrac{\pi}{6} \cos \alpha -\cos \dfrac{\pi}{6}\sin \alpha + \sin \dfrac{\pi}{6} \cos \alpha +\cos \dfrac{\pi}{6}\sin \alpha -\cos \alpha \\
				  & =\dfrac{1}{2}\cos \alpha + \dfrac{1}{2}\cos \alpha - \cos \alpha                                                                                \\
				  & =0.
			\end{array}$\\
		Vậy giá trị của biểu thức $P$ không phụ thuộc vào $\alpha$.

	}
\end{vd}
\begin{vd}%[DCHT Toán 11 - KNTT -Nguyễn Thành Nhân]%[1K1B2-2]
	Cho $\cos a=-\dfrac{2}{3}$ với $\dfrac{\pi}{2}<a<\pi$. Biết $S=\cos 2a+\sin 2a=m+n\sqrt{5}$ với $m$, $n\in \mathbb{Q}$ và $\dfrac{m}{n}=\dfrac{p}{q}$ là phân số tối giản. Tính $p-q$.
	\loigiai{
		Vì $\dfrac{\pi}{2}<a<\pi$ nên $\sin a>0$. Do đó
		\[\sin a=\sqrt{1-\cos^2 a}=\sqrt{1-\dfrac{4}{9}}=\dfrac{\sqrt{5}}{3}.\]
		Do đó
		\begin{eqnarray*}
			S&=& \cos^2 a-\sin^2 a+2\sin a\cdot \cos a\\
			&=& \dfrac{4}{9}-\dfrac{5}{9}+2\cdot \dfrac{-2}{3}\cdot \dfrac{\sqrt{5}}{3}\\
			&=&-\dfrac{1}{9}-\dfrac{4}{9}\cdot \sqrt{5}.
		\end{eqnarray*}
		Suy ra $m=-\dfrac{1}{9}$ và $n=-\dfrac{4}{9}$. Nên $\dfrac{m}{n}=\dfrac{1}{4}$. Suy ra $p=1$, $q=4$.\\
		Vậy $p-q=-3$.
	}
\end{vd}
\begin{vd}%[DCHT Toán 11 - KNTT -Nguyễn Thành Nhân]%[1K1B2-2]
	Cho $\cos a=\dfrac{5}{13}$ với $0<a<\dfrac{\pi}{2}$. Tính $\sin 2a$, $\cos 2a$, $\tan 2a$, $\sin \left(2a+\dfrac{\pi}{3}\right)$, $\tan\left(2a-\dfrac{\pi}{6}\right)$.
	\loigiai{
		Ta có $\sin^2 a=1-\cos^2 a=1-\dfrac{25}{169}=\dfrac{144}{169}\Leftrightarrow \hoac{&\sin a=\dfrac{12}{13}\\&\sin a=\dfrac{-12}{13}.}$\\
		Vì $0<a<\dfrac{\pi}{2}$ nên $\sin a>0$, suy ra $\sin a=\dfrac{12}{13}$.\\
		Từ đó ta có
		\begin{eqnarray*}
			\sin 2a &=& 2\sin a\cos a=2\cdot \dfrac{5}{13}\cdot \dfrac{12}{13}=\dfrac{120}{169};\\
			\cos 2a &=& 2\cos ^2 a-1=2\cdot \dfrac{25}{169}-1=-\dfrac{119}{169};\\
			\tan 2a &=& \dfrac{\sin 2a}{\cos 2a}=-\dfrac{120}{119};\\
			\sin \left(2a+\dfrac{\pi}{3}\right)&=&\sin 2a\cdot \cos \dfrac{\pi}{3}+ \cos 2a\cdot \sin \dfrac{\pi}{3}\\
			&=&\dfrac{120}{169}\cdot \dfrac{1}{2}-\dfrac{120}{119}\cdot \dfrac{\sqrt{3}}{2}=\dfrac{60(1-\sqrt{3})}{119}.\\
			\tan\left(2a-\dfrac{\pi}{6}\right)&=&\dfrac{\tan 2a-\tan \dfrac{\pi}{6}}{1+\tan 2a\tan \dfrac{\pi}{6}}=\dfrac{-\dfrac{120}{119}-\dfrac{\sqrt{3}}{3}}{1-\dfrac{120}{119}\cdot \dfrac{\sqrt{3}}{3}}\\
			&=&-\dfrac{360+119\sqrt{3}}{357-120\sqrt{3}}.
		\end{eqnarray*}
	}
\end{vd}
\begin{vd} %[DCHT Toán 11 - KNTT -Nguyễn Thành Nhân]%[1K1B2-2]
	Cho $\sin 2 a=\dfrac{3}{5}$ với $\dfrac{\pi}{2}<a<\pi$. Tính $\tan a+\cot a$, $\tan a-\cot a$.
	\loigiai{
		Ta có
		\begin{eqnarray*}
			\tan a+\cot a&=&\dfrac{\sin a}{\cos a}+\dfrac{\cos a}{\sin a}=\dfrac{\sin^2 a+\cos^2 a}{\sin a\cos a} \\
			&=&\dfrac{1}{\dfrac{1}{2}\sin 2a}=\dfrac{2}{\sin 2a}.
		\end{eqnarray*}
		Thay $\sin 2a=\dfrac{3}{5}$, ta được $\tan a+\cot a=\dfrac{2}{\dfrac{3}{5}}=\dfrac{10}{3}$.\\
		Ta cũng có
		\begin{eqnarray*}
			\tan a-\cot a&=&\dfrac{\sin a}{\cos a}-\dfrac{\cos a}{\sin a}=\dfrac{\sin^2 a-\cos^2 a}{\sin a\cos a} \\
			&=&\dfrac{-\cos 2a}{\dfrac{1}{2}\sin 2a}=-2\cot 2a.
		\end{eqnarray*}
		Vì $\sin 2a=\dfrac{3}{5}$ và $\dfrac{\pi}{2}<a<\pi$ nên $\cos 2a<0$, suy ra $\cos 2a=-\sqrt{1-\sin^2 2a} =-\dfrac{4}{5}$. \\
		Do đó $\tan a-\cot a=-\cot 2a=-\dfrac{-\dfrac{4}{5}}{\dfrac{3}{5}}=\dfrac{4}{3}$.
	}
\end{vd}
\begin{vd}%[DCHT Toán 11 - KNTT -Nguyễn Thành Nhân]%[1K1B2-2]
	Cho $\sin a+\cos a=m,\,(-\sqrt{2}\le m \le \sqrt{2})$. Tính $\left|\sin a-\cos a\right|$.
	\loigiai{
		Ta có $m^2=\left(\sin a+\cos a\right)^2=1+2\sin a\cos a=1+\sin 2a$. Suy ra $\sin 2a=m^2-1$. Do đó
		\begin{eqnarray*}
			\left|\sin a-\cos a\right|&=& \sqrt{\left(\sin a-\cos a\right)^2}=\sqrt{1-\sin 2a}=\sqrt{2-m^2}.
		\end{eqnarray*}

	}
\end{vd}
\begin{vd}
	\immini[thm]{Một sợi cáp $R$ được gắn vào một cột thẳng đứng ở vị trí cách mặt đất $14 \mathrm{~m}$. Một sợi cáp $S$ khác cũng được gắn vào cột đó ở vị trí cách mặt đất $12 \mathrm{~m}$. Biết rằng hai sợi cáp trên cùng được gắn với mặt đất tại một vị trí cách chân cột $15 \mathrm{~m}$ (Hình bên).
		\begin{tasks}
			\task Tính $\tan \alpha$, ở đó $\alpha$ là góc giữa hai sợi cáp trên.
			\task Tìm góc $\alpha$ (làm tròn kết quả đến hàng đơn vị theo đơn vị độ).
		\end{tasks}
		\cham{4}}{
		\begin{tikzpicture}[scale=0.8, font=\footnotesize,>=stealth]
			\path
			(5,0) coordinate (O)
			(0,0) coordinate (H)
			(0,4) coordinate (B)
			(0,6) coordinate (A)
			(1.5,1) coordinate (H1)
			(1.5,5) coordinate (B1)
			(1.5,7) coordinate (A1)
			;
			\draw[line width=4pt,cyan] (H)--(A) (H1)--(A1);
			\draw[red!70!black,line width=0.7pt] (O)--(A)node[pos=.5,above right]{$R$} (O)--(B)node[pos=.5,below]{$S$};
			\draw
			(A)--(A1) (B)--(B1) (H)--(H1)
			pic[draw,double,angle radius=11mm]{angle=A--O--H}
			pic[draw,angle radius=9mm]{angle=A--O--B};
			\draw[<->] (O)--(H)node[pos=.5,below]{$15$ m};
			\draw[<-] ($(H)+(-0.5,0)$)--($(0,1.5)+(-0.5,0)$)node[above]{\scriptsize$12$ m};
			\draw[->] ($(0,2)+(-0.5,0)$)--($(B)+(-0.5,0)$);
			\draw[<-] ($(H)+(-1,0)$)--($(0,2.5)+(-1,0)$)node[above]{\scriptsize$14$ m};
			\draw[->] ($(0,3)+(-1,0)$)--($(A)+(-1,0)$);
			\foreach \x/\g in {	O/-90,H/-90,B/150,A/90}\draw[fill=black] (\x) circle (.05) +(\g:.35)node{\footnotesize$\x$};
			\draw[dashed] (0,0)--(-1.5,0);
			\draw[dashed] (0,6)--(-1.5,6);
			\draw[dashed] (0,4)--(-0.6,4);
			\node[left] at (3.7,0.6) {$\beta$};
			\node[left] at (5.1,1.4) {$\alpha$};
			\draw[->] (4.6,1.2)--(4.25,0.65);
		\end{tikzpicture}}
	\loigiai{
		\begin{enumEX}[a)]{1}
			\item Đặt $\widehat{HOA}=\beta$ và $\widehat{HOB}=\beta_1$. Ta tính được $$\tan\beta = \dfrac{AH}{HO}=\dfrac{14}{15};\quad \tan\beta_1 = \dfrac{BH}{HO}=\dfrac{12}{15}.$$
			Khi đó
			$$\tan\alpha = \tan\left(\beta - \beta_1 \right)=\dfrac{\tan\beta- \tan\beta_1}{1+\tan\beta \tan\beta_1}=\dfrac{10}{131}.$$
			\item Với $\tan \alpha= \dfrac{10}{131}$ và góc $0<\alpha<90^\circ$, suy ra $\alpha= 4^\circ$.
		\end{enumEX}}
\end{vd}%\dongcham{4}
\begin{dang}{Sử dụng công thức biến đổi tích thành tổng}
\end{dang}

\begin{vd}
	Hãy tính giá trị của các biểu thức sau:
	\begin{enumEX}[a)]{1}
		\item $A=\cos 45^{\circ} \cos15^{\circ}$ \dotfill.
		\item $B=\cos 75^{\circ} \sin 15^{\circ}$ \dotfill.
		\item $C=\sin 75^{\circ} \sin 15^{\circ}$ \dotfill.
		\item $D=\sin \dfrac{11\pi}{12} \cos \dfrac{5\pi}{12}$ \dotfill.
	\end{enumEX}
	\loigiai{
		\begin{enumerate}[a)]
			\item $A=\dfrac{1}{2}\left(\cos(45^\circ+15^\circ)+\cos(45^\circ-15^\circ) \right)=\dfrac{1}{2}\left(\cos60^\circ+\cos30^\circ \right)=\dfrac{1+\sqrt{3}}{4}. $
			\item $B=\dfrac{1}{2}\left[\sin \left(75^{\circ}+15^{\circ}\right) -\sin \left(75^{\circ}-15^\circ\right)\right]=\dfrac{1}{2}\left(\sin90^\circ-\sin60^\circ \right)=\dfrac{2-\sqrt{3}}{4}.$
			\item $C=-\dfrac{1}{2}\left(\cos(75^\circ+15^\circ)-\cos(75^\circ-15^\circ) \right)=-\dfrac{1}{2}\left(\cos90^\circ-\cos60^\circ \right)=\dfrac{1}{4}$.
			\item $D=\sin \dfrac{11\pi}{12} \cos \dfrac{5\pi}{12}=\dfrac{1}{2}\left[\sin \left(\dfrac{16\pi}{12}\right) +\sin \left(\dfrac{6\pi}{12}\right)\right]= \dfrac{1}{2}\left[\sin \left(\dfrac{-2\pi}{3}\right) +\sin \left(\dfrac{\pi}{2}\right) \right]=  \dfrac{2- \sqrt{3}}{4}$.
		\end{enumerate}}
\end{vd} %\dongcham{14}

\begin{vd}%[0D6B3-3]%
	Biến đổi các biểu thức sau đây thành một tổng:
	\begin{listEX}[1]
		\item $\cos 5a\sin 3a$. \dotfill
		\item $2\cos (a+b)\cos (a-b)$.\dotfill
		\item $\sin (a-b)\cos (b-a)$.\dotfill
		\item $2\sin (a+b)\sin (a-b)$.\dotfill
	\end{listEX}
	\loigiai{
		\begin{listEX}[1]
			\item $\cos 5a\sin 3a=\sin 3a\cos 5a=\dfrac{1}{2}\left[\sin \left(3a-5a\right)+\sin \left(3a+5a\right)\right]=\dfrac{1}{2}\left[\sin \left(-2a\right)+\sin 8a\right]$.
			\item $2\cos (a+b)\cos (a-b)=2\cdot \dfrac{1}{2}\left[\cos \left(a+b-a+b\right)+\cos \left(a+b+a-b\right)\right]=\cos 2b+\cos 2a$.
			\item $\sin (a-b)\cos (b-a)=\dfrac{1}{2}\left[\sin (a-b-b+a)+\sin (a-b+b-a)\right]=\dfrac{1}{2}\sin (2a-2b)$.
			\item $2\sin (a+b)\sin (a-b)=2\cdot \dfrac{1}{2}\left[\cos \left(a+b-a+b\right)-\cos \left(a+b+a-b\right)\right]=\cos 2b-\cos 2a$.\\
			Vậy $2\sin (a+b)\sin (a-b)=\cos 2b-\cos 2a$.
		\end{listEX}
	}
\end{vd} %\dongcham{15}

\begin{vd}%%[DCHT Toán 11 - KNTT -Nguyễn Thành Nhân]%[1K1B2-3]
	Rút gọn các biểu thức sau
	\begin{listEX}[2]
		\item $A=\cos 11x\cos 3x-\cos 17x\cos 9x$.
		\item $B=\sin 18x\cos 3x-\sin 19x\cos 4x$.
		\item $C=\sin x\sin 3x+\sin 4x\sin 8x$.
		\item $D=\sin 2x\sin 6x-\cos x\cos 3x$.
		\item $E=\sin x\sin\left(\dfrac{\pi}{3}-x\right)\sin\left(\dfrac{\pi}{3}+x\right)$.
		\item $F=\cos \dfrac{x}{2}\cos \dfrac{3x}{2}-\sin x\sin 3x-\sin 2x\sin 3x$.
	\end{listEX}
	\loigiai{
		\begin{enumerate}
			\item $A=\dfrac{1}{2}\left( \cos 14x+\cos 8x \right) -\dfrac{1}{2} \left( \cos 26x+\cos 8x \right) =\dfrac{1}{2}\left(\cos 14x-\cos 26x\right)=\dfrac{1}{2}(-2)\sin 20x\sin (-6x) = \sin 20x \sin 6x$.
			\item $B=\dfrac{1}{2}\left( \sin 21x+\sin 15x \right) -\dfrac{1}{2} \left( \sin 23x+\sin 15x \right) =\dfrac{1}{2}\left(\sin 21x-\sin 23x\right)=-\cos 22x\sin x$.
			\item $C=\dfrac{1}{2} \left( \cos 2x-\cos 4x \right) +\dfrac{1}{2} \left( \cos 4x-\cos 12x \right) =-\dfrac{1}{2}\left(\cos 12x-\cos 2x\right)=\sin 7x\sin 5x$.
			\item $D=\dfrac{1}{2} \left( \cos 4x-\cos 8x\right) -\dfrac{1}{2}\left( \cos 4x+\cos 2x \right) =-\dfrac{1}{2}\left(\cos 8x+\cos 2x\right)=-\cos 5x\cos 3x$.
			\item $E=\sin x\cdot\dfrac{1}{2}\left(\cos (-2x) - \cos \dfrac{2\pi}{3}\right)= \dfrac{1}{2}\cos2x \sin x  + \dfrac{1}{4} \sin x =\dfrac{1}{4}\sin 3x-\dfrac{1}{4}\sin x+\dfrac{1}{4}\sin x=\dfrac{1}{4}\sin 3x$.
			\item $F=\dfrac{1}{2} \left( \cos 2x+\cos x \right)  - \dfrac{1}{2} \left( \cos 2x-\cos 4x \right) -\dfrac{1}{2} \left( \cos x-\cos 5x\right) =\dfrac{1}{2}\left(\cos 4x+\cos 5x\right)=\cos \dfrac{9x}{2}\cos \dfrac{x}{2}$.
		\end{enumerate}
	}
\end{vd}
\begin{dang}{Sử dụng công thức biến đổi tổng thành tích}
\end{dang}
\viduminhhoa
\begin{vd}%[0D6K3-3]%
	Chứng minh
	\begin{listEX}[1]
		\item $\sin 65^\circ +\sin 55^\circ= \sqrt{3}\cos 5^\circ$.
		\item $\cos 12^\circ -\cos 48^\circ=\sin 18^\circ$.
		\item $\sin 20^\circ -\sin 100^\circ+\sin 140^\circ = 0$.
		\item $\tan 9^\circ -\tan 27^\circ-\tan 63^\circ+\tan 81^\circ= 4$.
	\end{listEX}

	\loigiai{
		\begin{listEX}[1]
			\item Ta có $\sin 65^\circ +\sin 55^\circ=2 \cdot \sin 60^\circ \cdot \cos 5^\circ= \sqrt{3}\cos 5^\circ$.
			\item Ta có
			$ \cos 12^\circ -\cos 48^\circ=(-2) \cdot \sin 30^\circ \cdot \sin (-18^\circ) =\sin 18^\circ$.
			\item Ta có
			$\sin 20^\circ -\sin 100^\circ+\sin 140^\circ
				= \left( \sin 20^\circ +\sin 140^\circ \right) -\sin 100^\circ =2 \cdot \sin 80^\circ \cdot \cos 60^\circ-\sin 100^\circ =\sin 80^\circ -\sin 80^\circ = 0$.
			\item Với $\forall a \ne k90^\circ, k \in \mathbb{Z}$, ta có
			$$\tan a + \tan (90^\circ) = \tan a + \cot a = \dfrac{\sin a}{\cos a} + \dfrac{\cos a}{\sin a} = \dfrac{1}{\sin a \cos a} = \dfrac{2}{\sin 2a}. $$
			Suy ra
			\allowdisplaybreaks
			\begin{eqnarray*}
				\tan 9^\circ -\tan 27^\circ-\tan 63^\circ+\tan 81^\circ
				&=& \left( \tan 9^\circ +\tan 81^\circ\right) - \left( \tan 27^\circ+\tan 63^\circ\right)  \\
				&=& \dfrac{2}{\sin 18^\circ} - \dfrac{2}{\sin 54^\circ} \\
				&=& 2 \cdot \dfrac{\sin 54^\circ - \sin 18^\circ}{\sin 18^\circ \sin 54^\circ} \\
				&=& 2 \cdot \dfrac{2 \cos 36^\circ \sin 18^\circ}{\sin 18^\circ \sin 54^\circ} \\
				&=& 4\ (\text{ vì } 36^\circ + 54^\circ = 90^\circ \text{ nên } \cos 36^\circ = \sin 54^\circ).
			\end{eqnarray*}
		\end{listEX}
	}
\end{vd} \dongcham{23}
\begin{vd}%%[DCHT Toán 11 - KNTT -Nguyễn Thành Nhân]%[1K1B2-3]
	Biến đổi các tổng sau thành tích
	\begin{listEX}[2]
		\item $A=\sin 5x+\sin 6x+\sin 7x+\sin 8x$.
		\item $B=\sin x-\sin 3x+\sin 7x-\sin 5x$.
		\item $C=\cos 7x+\sin 3x+\sin 2x-\cos 3x$.
		\item $D=\sin 35^\circ +\cos 40^\circ +\sin 55^\circ +\cos 20^\circ$.
	\end{listEX}
	\loigiai{
		\begin{enumerate}
			\item $A=\left(\sin 8x+\sin 5x\right)+\left(\sin 7x+\sin 6x\right)=2\sin \dfrac{13x}{2}\cos \dfrac{3x}{2}+2\sin \dfrac{13x}{2}\cos \dfrac{x}{2}=2\sin \dfrac{13x}{2}\left(\cos \dfrac{3x}{2}+\cos \dfrac{x}{2}\right)$.
			\item $B=\left(\sin 7x+\sin x\right)-\left(\sin 5x+\sin 3x\right)=2\sin 4x\cos 3x-2\sin 4x\cos x=2\sin 4x\left(\cos 3x-\cos x\right)$.
			\item $C=\left(\cos 7x-\cos 3x\right)+\left(\sin 3x+\sin 2x\right)=-2\sin 5x\sin 2x+2\sin \dfrac{5x}{2}\cos \dfrac{x}{2}=2\sin \dfrac{5x}{2}\left(-2\cos \dfrac{5x}{2}\sin 2x+\cos \dfrac{x}{2}\right)$.
			\item $D=\left(\sin 55^\circ +\sin 35^\circ\right)+\left(\cos 40^\circ +\cos 20^\circ\right)=2\sin 45^\circ\cos 10^\circ +2\cos 30^\circ\cos 10^\circ =2\cos 10^\circ\left(\sin 45^\circ +\cos 30^\circ\right)$.
		\end{enumerate}
	}
\end{vd}
\begin{vd}
	Biến đổi các biểu thức sau đây thành một tích.
	\begin{tasks}(2)
		\task $A=\sin a+\sin 3a +\sin 5a$.
		\task $B=1+\cos x+\cos 2x+\cos 3x$.
	\end{tasks}
	\loigiai{
		\begin{enumerate}
			\item $\sin a+\sin 3a +\sin 5a=\sin 5a+\sin a +\sin 3a=2\sin 3a\cos 2a + \sin 3a = \sin 3a(2\cos 2a+1)$.\\
			      Vậy $A=\sin a+\sin 3a +\sin 5a = \sin 3a(2\cos 2a+1)$.
			\item $B=1+\cos x+\cos 2x+\cos 3x=\left(\cos 3x+\cos x\right)+\left(\cos 2x+1\right)$ \\
			      $=2\cos 2x\cos x+2\cos^2x-1+1=2\cos x\left(\cos 2x+\cos x\right)=2\cos x\cdot2\cos \dfrac{3x}{2}\cos\dfrac{x}{2}$ \\
			      Vậy $B=1+\cos x+\cos 2x+\cos 3x=4\cos x\cos \dfrac{3x}{2}\cos\dfrac{x}{2}$.
		\end{enumerate}
	}
\end{vd} \dongcham{10}

\begin{vd}%%[DCHT Toán 11 - KNTT -Nguyễn Thành Nhân]%[1K1B2-3]
	Rút gọn biểu thức $S=2\sin x\left(\cos x+\cos 3x+\cos 5x\right)$. Từ đó tính giá trị biểu thức $$P=\cos \dfrac{\pi}{7}+\cos \dfrac{3\pi}{7}+\cos \dfrac{5\pi}{7}.$$
	\loigiai{
		Ta có 
		\begin{eqnarray*}
			S&=&2\sin x\left(\cos x+\cos 3x+\cos 5x\right)\\
			&=& 2\sin x\cdot \cos x+2\sin x\cdot \cos 3x+2\sin x\cdot \cos 5x\\
			&=& \sin 2x+\sin 4x-\sin 2x+\sin 6x-\sin 4x\\
			&=& \sin 6x.
		\end{eqnarray*}
		Vậy $S=\sin 6x$.\\
		Áp dụng kết quả trên để tính $P=\cos \dfrac{\pi}{7}+\cos \dfrac{3\pi}{7}+\cos \dfrac{5\pi}{7}$.\\
		Vì $\sin \dfrac{\pi}{7}\ne 0$ nên
		\begin{eqnarray*}
			P&=&\cos \dfrac{\pi}{7}+\cos \dfrac{3\pi}{7}+\cos \dfrac{5\pi}{7}\\
			\Rightarrow  2P\cdot \sin \dfrac{\pi}{7}&=&2 \sin \dfrac{\pi}{7}\left(\cos \dfrac{\pi}{7}+\cos \dfrac{3\pi}{7}+\cos \dfrac{5\pi}{7}\right)\\
			&=& \sin \dfrac{6\pi}{7}=\sin \left(\pi-\dfrac{\pi}{7}\right)=\sin \dfrac{\pi}{7}\\
			\Rightarrow P&=& \dfrac{1}{2}.
		\end{eqnarray*}
		Vậy $P=\dfrac{1}{2}$.
	}
\end{vd}
\begin{vd}
	Một thiết bị trễ kĩ thuật số lặp lại tín hiệu đầu vào bằng cách lặp lại tín hiệu đó trong một khoảng thời gian cố định sau khi nhận được tín hiệu. Nếu một thiết bị như vậy nhận được nốt thuần $f_1(t)=5 \sin t$ và phát lại được nốt thuần $f_2(t)=5 \cos t$ thì âm kết hợp là $f(t)=f_1(t)+f_2(t)$, trong đó $t$ là biến thời gian. Chứng tỏ rằng âm kết hợp viết được dưới dạng $f(t)=k \sin (t+\varphi)$, tức là âm kết hợp là một sóng âm hình sin. Hãy xác định biên độ âm $k$ và pha ban đầu $\varphi(-\pi \leq \varphi \leq \pi)$ của sóng âm.
	\loigiai{
		Ta có $f(t)=f_1(t)+f_2(t)=5\left( \sin t + \cos t\right)=5\sqrt{2}\sin\left(t+\dfrac{\pi}{4} \right) $. Suy ra
		\begin{itemize}
			\item [$\bullet$] Biên độ âm $k=5\sqrt{2}$.
			\item [$\bullet$] Pha ban đầu $\varphi = \dfrac{\pi}{4}$.
		\end{itemize}
	}
\end{vd} \dongcham{8}

\begin{vd}
	Trong Vật lí, phương trình tổng quát của một vật dao động điều hoà cho bởi công thức $x(t)=A \cos (\omega t+\varphi)$, trong đó $t$ là thời điểm (tính bằng giây), $x(t)$ là li độ của vật tại thời điểm $t, A$ là biên độ dao động $(A>0)$ và $\varphi \in[-\pi ; \pi]$ là pha ban đầu của dao động. Xét hai dao động điều hoà có phương trình:
	$$
		\begin{aligned}
			 & x_1(t)=2 \cos \left(\dfrac{\pi}{3} t+\dfrac{\pi}{6}\right)(\mathrm{cm}),  \\
			 & x_2(t)=2 \cos \left(\dfrac{\pi}{3} t-\dfrac{\pi}{3}\right)(\mathrm{cm}) .
		\end{aligned}
	$$
	Tìm dao động tổng hợp $x(t)=x_1(t)+x_2(t)$ và sử dụng công thức biến đổi tổng thành tích để tìm biên độ và pha ban đầu của dao động tồng hợp này.
	\loigiai{
		Ta có $x(t)=x_1(t)+x_2(t)=2\left[ \cos \left(\dfrac{\pi}{3} t+\dfrac{\pi}{6}\right)+ \cos \left(\dfrac{\pi}{3} t-\dfrac{\pi}{3}\right)\right]=2\sqrt{2}\cos\left(\dfrac{\pi t}{3}-\dfrac{\pi}{12} \right) $. Suy ra
		\begin{itemize}
			\item [$\bullet$] Biên độ âm $k=2\sqrt{2}$.
			\item [$\bullet$] Pha ban đầu $\varphi = -\dfrac{\pi}{12}$.
		\end{itemize}
	}
\end{vd} \dongcham{8}

\subsection{BÀI TẬP TỰ LUẬN}

\begin{bt}
	Tính
	\begin{tasks}(1)
		\task $\cos\left(\alpha+\dfrac{\pi}{3}\right)$, biết $\sin \alpha=\dfrac{1}{\sqrt{3}}$ và $0<\alpha<\dfrac{\pi}{2}$.
		\task $\tan\left(\alpha -\dfrac{\pi}{4} \right)$, biết $\cos \alpha=-\dfrac{1}{3}$ và $\dfrac{\pi}{2}<\alpha<\pi$.
		\task $\cos\left(a+b\right),\sin\left(a-b\right)$, biết $\sin a=\dfrac{4}{5},0^\circ<a<90^\circ$ và $\sin b =\dfrac{2}{3},90^\circ<b<180^\circ$.
	\end{tasks}
	\loigiai{
		\begin{enumerate}[a)]
			\item Do $0<\alpha<\dfrac{\pi}{2}$ nên $\cos\alpha>0$.
			      Do đó $\cos\alpha=\sqrt{1-\sin^2\alpha}=\sqrt{\dfrac{2}{3}}=\dfrac{\sqrt{6}}{3}$.\\
			      Ta có $\cos\left(\alpha+\dfrac{\pi}{3}\right)=\cos\alpha\cos\dfrac{\pi}{3}-\sin\alpha\sin\dfrac{\pi}{3}=\dfrac{\sqrt{6}}{3}\cdot\dfrac{1}{2}-\dfrac{1}{\sqrt{3}}\cdot\dfrac{\sqrt{3}}{2}=\dfrac{-3+\sqrt{6}}{6}$.
			\item Do $\dfrac{\pi}{2}<\alpha<\pi$ nên $\sin\alpha>0$. Do đó $\sin\alpha=\sqrt{1-\cos^2\alpha}=\dfrac{2\sqrt{2}}{3}$. Suy ra $\tan\alpha=-2\sqrt{2}$.\\
			      Ta có $\tan\left(\alpha -\dfrac{\pi}{4} \right)=\dfrac{\tan\alpha-\tan\dfrac{\pi}{4}}{1+\tan\alpha\tan\dfrac{\pi}{4}}=\dfrac{\tan\alpha-1}{\tan\alpha+1}=\dfrac{-2\sqrt{2}-1}{-2\sqrt{2}+1}=\dfrac{9+4\sqrt{2}}{7}$.
			\item Có $0^\circ<a<90^\circ$ nên $\cos a=\dfrac{3}{5}$, $90^\circ<b<180^\circ$ nên $\cos b=-\dfrac{\sqrt{5}}{3}$.\\
			      $\cos\left(a+b\right)=\cos a\cos b -\sin a\sin b = \dfrac{3}{5}\cdot\dfrac{-\sqrt{5}}{3}-\dfrac{4}{5}\cdot\dfrac{2}{3}=-\dfrac{8+3\sqrt{5}}{15}$.\\
			      $\sin\left(a-b\right)=\sin a\cos b-\sin b\cos a=\dfrac{4}{5}\cdot\dfrac{-\sqrt{5}}{3}-\dfrac{2}{3}\cdot\dfrac{3}{5}=-\dfrac{6+4\sqrt{5}}{15}$.
		\end{enumerate}
	}
\end{bt} \cham{19}


\begin{bt}
	Chứng minh các đẳng thức sau
	\begin{enumEX}{2}
		\item $\sin^4\alpha +\cos^4\alpha =1-\dfrac{1}{2}\sin^22\alpha$;
		\item $\sin^6\alpha +\cos^6\alpha =\dfrac{5}{8}+\dfrac{3}{8}\cos 4\alpha $.
	\end{enumEX}
	\loigiai{
		\begin{enumerate}[a)]
			\item Ta có $\sin^4\alpha +\cos^4\alpha =\left(\sin^2\alpha +\cos^2\alpha \right)^2-2\sin^2\alpha \cos^2\alpha =1-\dfrac{1}{2}\sin^22\alpha$.
			\item Ta có $\sin^6\alpha +\cos^6\alpha =\left(\sin^2\alpha \right)^3+\left(\cos^2\alpha \right)^3+3\sin^2\alpha \cos^2\alpha \left(\sin^2\alpha +\cos^2\alpha \right)$\\
			      \phantom{Ta có $\sin^6\alpha +\cos^6\alpha =\left(\sin^2\alpha \right)^3+\left(\cos^2\alpha \right)^3$} $-3\sin^2\alpha \cos^2\alpha \left(\sin^2\alpha +\cos^2\alpha \right)$ \\
			      \phantom{Ta có $\sin^6\alpha +\cos^6\alpha$} $=\left(\sin^2\alpha +\cos^2\alpha \right)^3-3\sin^2\alpha \cos^2\alpha =1-\dfrac{3}{4}\left(2\sin \alpha \cos \alpha \right)^2$\\
			      \phantom{Ta có $\sin^6\alpha +\cos^6\alpha$} $=1-\dfrac{3}{4}\sin^22\alpha =1-\dfrac{3}{8}\left(1-\cos 4\alpha \right)=\dfrac{5}{8}+\dfrac{3}{8}\cos 4\alpha$.
		\end{enumerate}
	}
\end{bt} \cham{15}

\begin{bt}
	Chứng minh các đẳng thức sau
	\begin{enumEX}{2}
		\item $2\sin \left(\dfrac{\pi}{4}+\alpha \right)\sin \left(\dfrac{\pi}{4}-\alpha \right)=\cos 2\alpha $;
		\item $\sin \alpha \left(1+\cos 2\alpha \right)=\sin 2\alpha \cos \alpha $;
		\item $\dfrac{1+\sin 2\alpha -\cos 2\alpha}{1+\sin 2\alpha +\cos 2\alpha}=\tan \alpha $;
		\item $\tan \alpha -\dfrac{1}{\tan \alpha}=-\dfrac{2}{\tan 2\alpha}$.
	\end{enumEX}
	\loigiai{
		\begin{enumerate}[a)]
			\item Ta có $2\sin \left(\dfrac{\pi}{4}+\alpha \right)\sin \left(\dfrac{\pi}{4}-\alpha \right)=2\left[\dfrac{\sqrt{2}}{2}\left(\cos \alpha +\sin \alpha \right)\dfrac{\sqrt{2}}{2}\left(\cos \alpha -\sin \alpha \right)\right]$\\
			      \phantom{Ta có $2\sin \left(\dfrac{\pi}{4}+\alpha \right)\sin \left(\dfrac{\pi}{4}-\alpha \right)$} $=\cos^2\alpha -\sin^2\alpha =\cos 2\alpha $.
			\item Ta có $\sin \alpha \left(1+\cos 2\alpha \right)=\sin \alpha \left(1+2\cos^2\alpha -1\right)=2\sin \alpha \cos^2\alpha =\sin 2\alpha \cos \alpha $.
			\item Ta có $\dfrac{1+\sin 2\alpha -\cos 2\alpha}{1+\sin 2\alpha +\cos 2\alpha}=\dfrac{\sin 2\alpha +\left(1-\cos 2\alpha \right)}{\sin 2\alpha +\left(1+\cos 2\alpha \right)}=\dfrac{\sin 2\alpha +2\sin^2\alpha}{\sin 2\alpha +2\cos^2\alpha}$\\
			      \phantom{Ta có $\dfrac{1+\sin 2\alpha -\cos 2\alpha}{1+\sin 2\alpha +\cos 2\alpha}$} $=\dfrac{2\sin \alpha \left(\cos \alpha +\sin \alpha \right)}{2\cos \alpha \left(\sin \alpha +\cos \alpha \right)}=\tan \alpha $.
			\item Ta có $\tan \alpha -\dfrac{1}{\tan \alpha}=2\cdot\dfrac{\tan^2\alpha -1}{2\cdot\tan \alpha}=-\dfrac{2}{\tan 2\alpha}$.
		\end{enumerate}
	}
\end{bt} \cham{30}
\begin{bt}%[VDT]%[DCHT Toán 11 - KNTT -Phong Lê] %[1K1K2-1] 
	Cho tam giác $ABC$ có $\cos B = \dfrac{3}{5}$, $\cos C = \dfrac{\sqrt{21}}{5}$. Chứng minh rằng $$\sin A =\sin B \cos C + \cos B \sin C$$ và tính $\sin A$.
	\dapso{$\sin A = \dfrac{6+4\sqrt{21}}{25}$.}
	\loigiai{
		Theo tính chất tổng ba góc trong một tam giác, ta có $A+B+C = \pi$.\\
		Do đó
		$\sin A = \sin (\pi - A) =\sin (B+C) = \sin B\cos C + \cos B \sin C.$	\\
		Vậy $\sin A =\sin B \cos C + \cos B \sin C$.\\
		Ta có $\sin^2B+\cos^2B = 1$, suy ra $\sin^2B = 1-\cos^2B = 1-\left(\dfrac{3}{5}\right)^2 = \dfrac{16}{25}$.\\
		Mà $0<B<\pi$ nên $\sin B >0$, suy ra $\sin B =\sqrt{\dfrac{16}{25}} = \dfrac{4}{5}$.\\
		Lại có $\sin^2C+\cos^2C = 1$, suy ra $\sin^2C = 1-\cos^2C = 1-\left(\dfrac{\sqrt{21}}{5}\right)^2 = \dfrac{4}{25}$.\\
		Mà $0<C<\pi$ nên $\sin C >0$, suy ra $\sin C =\sqrt{\dfrac{4}{25}} = \dfrac{2}{5}$.\\
		Do đó $\sin A = \sin B \cos C + \cos B \sin C = \dfrac{4}{5}.\dfrac{\sqrt{21}}{5} + \dfrac{3}{5}.\dfrac{2}{5} = \dfrac{6+4\sqrt{21}}{25}$.\\
		Vậy $\sin A = \dfrac{6+4\sqrt{21}}{25}$.
	}

\end{bt}

\begin{bt}%[DCHT Toán 11 - KNTT -Nguyễn Thành Nhân]%[1K1B2-2]
Rút gọn các biểu thức sau
\begin{listEX}[2]
\item $A=\sin x\cos x\cos 2x$.
\item $B=\cos^42x-\sin^42x$.
\item $C=4\sin x\cdot\sin\left(x+\dfrac{\pi}{2}\right)\cdot\sin\left(2x+\dfrac{\pi}{2}\right)$.
\item $D=\sin 2x+\cos 2x-2\cos x\left(\sin x+\cos x\right)+1$.
\end{listEX}
\loigiai{
\begin{enumerate}
\item $A=\dfrac{1}{2}\cdot 2\sin x\cos x\cos 2x=\dfrac{1}{2}\sin 2x\cos 2x=\dfrac{1}{4}\cdot 2\sin 2x\cos 2x=\dfrac{1}{4}\sin 4x$.
\item $B=\left(\cos^22x+\sin^22x\right)\left(\cos^22x-\sin^22x\right)=\cos 4x$.
\item $C=4\sin x\cdot\cos x\cdot\cos 2x=2\sin 2x\cdot\cos 2x=\sin 4x$.
\item $D=\sin 2x+\cos 2x-\sin 2x-2\cos^2x+1=\cos 2x-\cos 2x=0$.
\end{enumerate}
}
\end{bt}
\begin{bt}%[DCHT Toán 11 - KNTT -Nguyễn Thành Nhân]%[1K1B2-2]
Cho $\cos 2x=\dfrac{1}{3}$ Tính giá trị các biểu thức sau
\begin{listEX}[2]
\item $A=\sin^2x\cdot\cos^2x$.
\item $B=\dfrac{1+\sin^2x}{\cos^2x}$.
\item $C=\dfrac{1+\cot^2x}{1-\cot^2x}$.
\item $D=\sin^6x+\cos^6x$.
\end{listEX}
\loigiai{
\begin{enumerate}
\item $A=\dfrac{1}{2}\left(1-\cos 2x\right)\cdot\dfrac{1}{2}\left(1+\cos 2x\right)=\dfrac{1}{4}\left(1-\cos^22x\right)=\dfrac{1}{4}\left(1-\dfrac{1}{9}\right)=\dfrac{2}{9}$.
\item $B=\dfrac{1+\dfrac{1}{2}\left(1-\cos 2x\right)}{\dfrac{1}{2}\left(1+\cos 2x\right)}=\dfrac{3-\cos 2x}{1+\cos 2x}=\dfrac{3-\dfrac{1}{3}}{1+\dfrac{1}{3}}=2$.
\item $C=\dfrac{\sin^2x+\cos^2x}{\sin^2x-\cos^2x}=\dfrac{-1}{\cos 2x}=\dfrac{-1}{\dfrac{1}{3}}=-3$.
\item \begin{eqnarray*}
D&=&\left(\sin^2x+\cos^2x\right)^3-3\sin^2x\cos^2x\left(\sin^2x+\cos^2x\right)\\
&=&1-\dfrac{3}{4}\left(1-\cos^22x\right)=\dfrac{1}{4}+\dfrac{3}{4}\cos^22x\\
&=& \dfrac{1}{4}+\dfrac{3}{4}\cdot \dfrac{1}{9}=\dfrac{1}{3}.
\end{eqnarray*}
\end{enumerate}
}
\end{bt}
\begin{bt}%[DCHT Toán 11 - KNTT -Nguyễn Thành Nhân]%[1K1K2-2]
Chứng minh đẳng thức $\sin^6 x \cos^2 x+\sin^2 x\cos^6 x=\dfrac{1}{8}\left(1-\cos^4 2x\right)$.
\loigiai{Biến đổi vế trái, ta có
\begin{eqnarray*}
VT&=& \sin^6 x \cos^2 x+\sin^2 x\cos^6 x\\
&=& \sin^2 x\cos^2 x\left(\sin^4 x+\cos^4 x\right)\\
&=& \dfrac{1}{4}\cdot \sin^2 2x\left(1-2\sin^2 x\cos ^2 x\right)\\
&=& \dfrac{1}{4}\cdot \sin^2 2x\left(1-\dfrac{1}{2}\sin^2 2x\right)\\
&=& \dfrac{1}{8}\cdot \sin^2 2x\left(2-\sin^2 2x\right)\\
&=& \dfrac{1}{8}\cdot \left(1-\cos^2 2x\right)\cdot \left(1+\cos^2 2x\right)\\
&=&\dfrac{1}{8}\left(1-\cos^4 2x\right)=VP.
\end{eqnarray*}

}
\end{bt}

\begin{bt}
Tính giá trị các biểu thức:
\begin{enumEX}{2}
\item $A=\cos \dfrac{2\pi}{7}+\cos \dfrac{4\pi}{7}+\cos \dfrac{6\pi}{7}$;
\item $B=\tan 9^\circ-\tan 27^\circ-\tan 63^\circ+\tan 81^\circ$.
\end{enumEX}
\loigiai{
\begin{enumerate}
\item Ta có
$$ \begin{aligned}
A\cdot\sin \dfrac{\pi}{7} &=\sin \dfrac{\pi}{7} \cos \dfrac{2\pi}{7}+\sin \dfrac{\pi}{7} \cos \dfrac{4\pi}{7}+\sin \dfrac{\pi}{7} \cos \dfrac{6\pi}{7} \\
			&=\dfrac{1}{2}\left[\sin \dfrac{3\pi}{7}-\sin \dfrac{\pi}{7}+\sin \dfrac{5\pi}{7}-\sin \dfrac{3\pi}{7}+\sin \dfrac{7\pi}{7}-\sin \dfrac{5\pi}{7}\right] \\
			&=\dfrac{1}{2}\left(-\sin \dfrac{\pi}{7}\right)=-\dfrac{1}{2}\sin \dfrac{\pi}{7}\\
			\Rightarrow A&=-\dfrac{1}{2}.
\end{aligned}$$
\item Ta có
\begin{align*}
B&=\tan 9^\circ-\tan 27^\circ-\tan 63^\circ+\tan 81^\circ=\tan 9^\circ+\tan 81^\circ-\left(\tan 27^\circ+\tan 63^\circ\right)\cr
&=\dfrac{\sin 9^\circ}{\cos 9^\circ}+\dfrac{\sin 81^\circ}{\cos 81^\circ}-\left(\dfrac{\sin 27^\circ}{\cos 27^\circ}+\dfrac{\sin 63^\circ}{\cos 63^\circ}\right)\cr
&=\dfrac{\sin 9^\circ\cos 81^\circ+\cos 9^\circ\sin 81^\circ}{\cos 9^\circ\cos 81^\circ}-\dfrac{\sin 27^\circ\cos 63^\circ+\cos 27^\circ\sin 63^\circ}{\cos 27^\circ\cos 63^\circ}\cr
&=\dfrac{\sin 90^\circ}{\cos 9^\circ\sin 9^\circ}-\dfrac{\sin 90^\circ}{\cos 27^\circ\sin 27^\circ}=\dfrac{2}{\sin 18^\circ}-\dfrac{2}{\sin 54^\circ}=\dfrac{2\left(\sin 54^\circ-\sin 18^\circ\right)}{\sin 54^\circ\sin 18^\circ}\cr
&=\dfrac{4\cos 36^\circ\sin 18^\circ}{\sin 54^\circ\sin 18^\circ}=4 \left(\text{ do } \cos 36^\circ=\sin 54^\circ\right).
\end{align*}
\end{enumerate}
}
\end{bt}
\begin{bt}%[DCHT Toán 11 - KNTT -Nguyễn Thành Nhân]%[1K1B2-2]
Chứng minh các đẳng thức sau
\begin{listEX}[1]
\item $8\sin^4x=3-4\cos 2x+\cos 4x$.
\item $\sin 4x=4\sin x\cdot\cos x\left(1-2\sin^2x\right)$.
\end{listEX}
\loigiai{
\begin{enumerate}
\item $VP=3-4\left(1-2\sin^2x\right)+2\cos^22x-1=-2+8\sin^2x+2\left(1-2\sin^2x\right)^2=8\sin^4x=VT$.
\item $VP=2\sin 2x\cdot\cos 2x=\sin 4x=VT$.
\end{enumerate}
}
\end{bt}
\begin{bt}%[DCHT Toán 11 - KNTT -Nguyễn Thành Nhân]%[1K1B2-2]
Chứng minh đẳng thức $\dfrac{\sin x+\cos x-1}{\sin x-\cos x+1}=\dfrac{\cos x}{1+\sin x}$.
\loigiai{
Biến đổi vế trái, ta được
\begin{eqnarray*}
VT&=& \dfrac{\sin x+\cos x-1}{\sin x-\cos x+1}=\dfrac{2\sin \dfrac{x}{2}\cos \dfrac{x}{2}-2\sin^2 \dfrac{x}{2}}{2\sin \dfrac{x}{2}\cos \dfrac{x}{2}+2\sin^2 \dfrac{x}{2}}\\
&=& \dfrac{2\sin \dfrac{x}{2}\left(\cos \dfrac{x}{2}-\sin \dfrac{x}{2}\right)}{2\sin \dfrac{x}{2}\left(\cos \dfrac{x}{2}+\sin \dfrac{x}{2}\right)}=\dfrac{ \cos \dfrac{x}{2}-\sin \dfrac{x}{2}}{\cos \dfrac{x}{2}+\sin \dfrac{x}{2}}\\
&=& \dfrac{\cos^2 \dfrac{x}{2}-\sin^2 \dfrac{x}{2}}{\left(\cos \dfrac{x}{2}+\sin \dfrac{x}{2}\right)^2}=\dfrac{\cos x}{1+\sin x}=VP.
\end{eqnarray*}
}
\end{bt}
\begin{bt}%[DCHT Toán 11 - KNTT -Nguyễn Thành Nhân]%[1K1K2-2]
Chứng minh đẳng thức $\sin^2 \left(\dfrac{\pi}{8}+x\right)-\sin^2 \left(\dfrac{\pi}{8}-x\right)=\dfrac{\sqrt{2}}{2}\sin 2x$.
\loigiai{
Áp dụng công thức hạ bậc, biến đổi vế trái, ta được
\begin{eqnarray*}
VT&=& \sin^2 \left(\dfrac{\pi}{8}+x\right)-\sin^2 \left(\dfrac{\pi}{8}-x\right)\\
&=& \dfrac{1-\cos \left(\dfrac{\pi}{4}+2x\right)}{2}-\dfrac{1-\cos \left(\dfrac{\pi}{4}-2x\right)}{2}\\
&=& \dfrac{\cos \left(\dfrac{\pi}{4}-2x\right)-\cos \left(\dfrac{\pi}{4}+2x\right)}{2}\\
&=& \dfrac{-2 \sin \dfrac{\pi}{4}\sin (-2x)}{2}=\dfrac{\sqrt{2}}{2}\sin 2x=VP.
\end{eqnarray*}
}
\end{bt}
\begin{bt}%%[DCHT Toán 11 - KNTT -Nguyễn Thành Nhân]%[1C1K2-3] 
	Chứng minh rằng với mọi tam giác $ABC$ ta luôn có
	$$\sin A+\sin B-\sin C=4\sin \dfrac{A}{2}\sin \dfrac{B}{2}\cos \dfrac{C}{2}.$$
	\loigiai{\\
		Ta có $\sin A+\sin B-\sin C=2\sin \dfrac{A+B}{2} \cos \dfrac{A-B}{2}-\sin \left(2\cdot \dfrac{C}{2}\right)$ \\
		Vì $\dfrac{A+B}{2}+\dfrac{C}{2}=90^\circ$ nên $\sin \dfrac{A+B}{2}=\cos \dfrac{C}{2}$ và $\sin \dfrac{C}{2}=\cos \dfrac{A+B}{2}$.\\
		Từ đó $\sin A+\sin B-\sin C=2\cos\dfrac{C}{2}\cos\dfrac{A-B}{2}-2\sin\dfrac{C}{2}\cos\dfrac{C}{2}$\\
		$=2\cos\dfrac{C}{2}\left[\cos\dfrac{A-B}{2}-\cos\dfrac{A+B}{2} \right]$\\
		$=2\cos \dfrac{C}{2}\cdot (-2)\sin \dfrac{A}{2} \sin \left(-\dfrac{B}{2}\right)=4\sin \dfrac{A}{2} \sin \dfrac{B}{2} \cos \dfrac{C}{2}$.\\
		Vậy $\sin A+\sin B-\sin C=4\sin \dfrac{A}{2}\sin \dfrac{B}{2}\cos \dfrac{C}{2}$.
	}
\end{bt}
\begin{bt}%%[DCHT Toán 11 - KNTT -Nguyễn Thành Nhân]%[1K1K2-3] 
	Chứng minh rằng với mọi tam giác nhọn $ABC$ ta luôn có
	$$\dfrac{\sin A+\sin B-\sin C}{\cos A+\cos B-\cos C+1}=\tan \dfrac{A}{2} \tan \dfrac{B}{2} \cot \dfrac{C}{2}.$$
	\loigiai{\\
		Ta có $\dfrac{\sin A+\sin B-\sin C}{\cos A+\cos B-\cos C+1}=\dfrac{2\sin \dfrac{A+B}{2} \cos \dfrac{A-B}{2}-\sin \left(2\cdot \dfrac{C}{2}\right)}{2\cos \dfrac{A+B}{2} \cos \dfrac{A-B}{2}+1-\cos \left(2 \dfrac{C}{2}\right)}$ \\
		$\dfrac{2\cos \dfrac{C}{2} \cos \dfrac{A-B}{2}-2\sin \dfrac{C}{2} \cos \dfrac{C}{2}}{2\sin \dfrac{C}{2} \cos \dfrac{A-B}{2}+2\sin^2\dfrac{C}{2}}=\dfrac{2\cos \dfrac{C}{2} \left(\cos \dfrac{A-B}{2}-\cos \dfrac{A+B}{2}\right)}{2\sin \dfrac{C}{2} \left(\cos \dfrac{A-B}{2}+\cos \dfrac{A+B}{2}\right)}$ \\
		$=\cot \dfrac{C}{2} \dfrac{-2\sin \dfrac{A}{2} \sin \left(-\dfrac{B}{2}\right)}{2\cos \dfrac{A}{2} \cos \left(-\dfrac{B}{2}\right)}=\cot \dfrac{C}{2} \tan \dfrac{A}{2} \tan \dfrac{B}{2}=\tan \dfrac{A}{2} \tan \dfrac{B}{2} \cot \dfrac{C}{2}$.\\
		Vậy $\dfrac{\sin A+\sin B-\sin C}{\cos A+\cos B-\cos C+1}=\tan \dfrac{A}{2} \tan \dfrac{B}{2} \cot \dfrac{C}{2}$.
	}
\end{bt}
\begin{bt}
	\immini[thm]{Tam giác $ABC$ vuông tại $B$ và có hai cạnh góc vuông là $AB=4$, $BC=3$. Vẽ điểm $D$ nằm trên tia đối của tia $CB$ thỏa mãn $\widehat{CAD}=30^{\circ}$. Tính $\tan \widehat{BAD}$, từ đó tính độ dài cạnh $CD$ (\textit{làm tròn đến hàng phần chục}).}{\hspace{1cm}
		\begin{tikzpicture}[scale=0.6, font=\footnotesize,>=stealth]
			\path
			(0,0) coordinate (B)
			(2.5,0) coordinate (C)
			(6,0) coordinate (D)
			(0,4) coordinate (A)
			;
			\draw
			(C)--(A)--(D)--cycle
			(A)--(B)node[pos=.5,left]{$4$}
			(B)--(C)node[pos=.5,below]{$3$}
			pic[draw,angle radius=2mm]{right angle=A--B--D}
			pic[draw,angle radius=6mm]{angle=C--A--D};
			\foreach \x/\g in {B/-90,C/-90,D/-90,A/90}\draw[fill=black] (\x) circle (.05) +(\g:.5)node{\footnotesize$\x$};
		\end{tikzpicture}}
	\loigiai{
	Đặt $\widehat{BAC}=\varphi$ thì $\tan\varphi=\dfrac{BC}{AB}=\dfrac{3}{4}$. Theo hình vẽ
	$$\tan \widehat{BAD} =\tan\left(\varphi + 30^\circ \right)=\dfrac{\tan \varphi + \tan 30^\circ }{1-\tan \varphi \tan 30^\circ} =\dfrac{48+25\sqrt{3}}{39}.$$
	Ta có  $\tan \widehat{BAD}=\dfrac{BD}{AB} \Rightarrow BD=AB\tan \widehat{BAD}$. Khi đó
	$$CD=BD-BC=AB\tan \widehat{BAD}-3=6{,}4.$$
	}
\end{bt} \cham{8}

\begin{bt}
	Dao động của một vật là tổng hợp của hai dao động điều hòa cùng phương, có phương trình lần lượt là $x_1=6 \cos 100 \pi t$ (mm) và $x_2=6 \sin 100 \pi t$ (mm), ($t$ tính bằng giây). Tính li độ của vật tại thời điểm $\mathrm{t}=0,25$ giây.
	\dapso{$-6 \mathrm{~mm}$.}
	\loigiai{
		Phương trình dao động tổng hợp của vật là:
		$$ x = x_1 + x_2 = 6\cos(100\pi t) + 6\sin(100\pi t). $$
		Tại thời điểm $t = 0{,}25$ giây, li độ của vật là:
		$$ x = 6\cos(100\pi \cdot 0{,}25) + 6\sin(100\pi \cdot 0{,}25) $$
		$$ x = 6\cos(25\pi) + 6\sin(25\pi). $$
		Ta có $\cos(25\pi) = \cos(\pi) = -1$ và $\sin(25\pi) = \sin(\pi) = 0$.
		Do đó:
		$$ x = 6 \cdot (-1) + 6 \cdot 0 = -6 \text{ (mm)}. $$
		Vậy li độ của vật tại thời điểm $t = 0{,}25$ giây là $-6$ mm.}
\end{bt} \cham{8}


\begin{bt}%[VDC]%[DCHT Toán 11 - KNTT -Phong Lê] %[1K1G2-1]
	Một vật thực hiện đồng thời hai dao động điều hòa có phương trình $x_1(t) =2\sqrt{3} \sin \left(4\pi t+\dfrac{\pi}{6}\right)$ và $x_2(t) = 2\cos \left(4\pi t+\dfrac{\pi}{6}\right)$. Chứng tỏ rằng phương trình dao động tổng hợp của vật đó $x(t) = x_1(t)+x_2(t)$ viết được dưới dạng $x(t) = A\cos (\omega t + \varphi)$, tức là dao động tổng hợp của vật đó là dao động điều hòa. Hãy xác định biên độ $A$, tần số góc $\omega$ và pha ban đầu $\varphi$ ($-\pi<\varphi<\pi$) của dao động tổng hợp.
	\dapso{$A=4, \omega = 4\pi, \varphi = -\dfrac{\pi}{6}$.}
	\loigiai{
		Ta có $x(t) = x_1(t)+x_2(t) = 2\sqrt{3}\sin \left(4\pi t+\dfrac{\pi}{6}\right) + 2\cos \left(4\pi t+\dfrac{\pi}{6}\right)$, đồng thời\\
		$\dfrac{1}{2}\cos\left(4\pi t+\dfrac{\pi}{6}\right) + \dfrac{\sqrt{3}}{2}\sin \left(4\pi t+\dfrac{\pi}{6}\right) =\cos \left(4\pi t+\dfrac{\pi}{6}\right) \cos \dfrac{\pi}{3} + \sin \left(4\pi t+\dfrac{\pi}{6}\right)\sin \dfrac{\pi}{3}  = \cos \left(4\pi t-\dfrac{\pi}{6}\right) $,\\
		suy ra $x(t) = 4\left[\dfrac{1}{2}\cos\left(4\pi t+\dfrac{\pi}{6}\right) + \dfrac{\sqrt{3}}{2}\sin \left(4\pi t+\dfrac{\pi}{6}\right)\right] = 4\cos\left(4\pi t - \dfrac{\pi}{6}\right).$\\
		Vậy dao động tổng hợp $x(t)$ có biên độ $A = 4$, tần số góc $\omega = 4\pi$ và pha ban đầu $\varphi = -\dfrac{\pi}{6}$.
	}
\end{bt}